\section{The energy method with prescribed columns}\label{sec:energy}

This section recalls OpenAI's energy method, extends it to determinants with prescribed columns, and
proves Theorem~\ref{thm:real} assuming the bound of Section~\ref{sec:signed}. Throughout, $\alpha_1=\AlphaOne$.

\subsection{OpenAI's energy}\label{ssec:energy}
By Andr\'eief's identity, $\det\cA_N$ is an integral over $n$ points $x_i\in(-1,1)$ and $n$ points
$s_j\in(0,1)$. OpenAI bounds the integrand by $e^{n^2E}$, where $E$ is a discrete logarithmic energy of the
points, written in the coordinate $x$ with $t=2x/(1+x^2)$ \cite[\S5]{OpenAI2026}. It has the form
\begin{multline*}
 n^2E=\sum_iV_1(x_i)+\sum_jV_2(s_j)+\sum_{i<i'}\bigl(\log|x_i-x_{i'}|+\log(1-x_ix_{i'})\bigr)\\
 +\sum_{j<j'}\log|s_j-s_{j'}|-\sum_{i,j}\log(1-2x_is_j+x_i^2),
\end{multline*}
schematically, with external fields $V_1,V_2$ of order $n$. The first two sums are the
external fields. The next two are the logarithmic interactions among points of the same kind, and the last
is the interaction between the two kinds. There are two cases, according to which of OpenAI's two bounds for
the determinant of the row functions is used \cite[Props.~5.2, 5.3]{OpenAI2026}. The case with doubled interaction is the one in which OpenAI's condition
$\sum_i(1-x_i^2)/(1+x_i^2)\le n-1-2g$ \cite[Eq.~(67)]{OpenAI2026} holds, and the constrained case is its complement. In the \emph{constrained
case} the row field acquires a term with a multiplier $\lambda_1$ for the constraint that defines the case.
In the \emph{case with doubled interaction} the interaction among the $x_i$ has weight $2$ instead of $1$. This case also bounds a determinant
$\det\mathsf L$ of size $n$ by $e^{O(n\log n)}$, which is the factor that Section~\ref{sec:signed} improves.

The energy is bounded from above by duality. Given trial sequences $u=(u_k)$ and $v=(v_k)$, one writes the
interaction terms as quadratic forms in the moments $\frac1n\sum_iT_k(x_i)$ and $\frac1n\sum_js_j^k$. Then one
completes the squares against $u$ and $v$. What remains is the sum of the norms $\|u\|^2=\sum_ku_k^2/k$ and
$\|v\|^2$ and, for each point, a function of one variable. With $\mathsf T(u,x)=\sum_ku_kT_k(x)/k$ and
$\mathsf S(v,x)=\sum_kv_kx^k/k$, these functions are, in the case with doubled interaction,
\begin{align*}
 X(x)&=\alpha_A\log|x|+\beta_h\log(1-x)-\gamma_H\log(1+x^2)-4\mathsf T(u,x)-\mathsf S(v,x),\\
 Y(s)&=\tfrac{b_0}{n_0}\log s+\tfrac{q_0}{n_0}\log(1-s)+2\mathsf T(v,s).
\end{align*}
In the constrained case, with a trial sequence $u'$ in place of $u$, they are
\begin{multline*}
 X_1(x)=\alpha_A\log|x|+\beta_h\log(1-x)-(\gamma_H-\tfrac12)\log(1+x^2)\\
 +\lambda_1\bigl(\tfrac{2g_0}{n_0}-\tfrac{2x^2}{1+x^2}\bigr)-2\mathsf T(u',x)-\mathsf S(v,x),
\end{multline*}
with the same $Y$. If $X\le M_X$ and $Y\le M_Y$ on the whole domains, the energy of every configuration is at
most a constant determined by $M_X$, $M_Y$ and the norms. In the language of potential theory the extremal
configuration is described by an equilibrium problem for two interacting measures in an external field
\cite{SaffTotik}. The best trial sequences are the moments of its solution, as Carreira observed
\cite{Carreira2026}. Any trial sequences give a valid upper bound.

\subsection{Prescribed points}\label{ssec:prescribed}
When $m$ of the $n$ columns are point values, only $n-m$ points of each kind are integrated over, and the
energy is still normalized by $n$. The prescribed points enter through the row determinant, now evaluated
at the free points together with the prescribed ones. There are two kinds of prescribed points: the
quadrature nodes, and points $y\in(0,1)$. In the coordinate $x$ the nodes lie in the set $E_W$ of the $x$ with
$|x|\le1$ and $2x/(1+x^2)$ in one of the intervals of Section~\ref{sec:real}. We have $E_W\subset E_x$. Let
$\alpha=m/n$ be the fraction of prescribed columns and $\alpha_2$ the fraction of prescribed points of the
second kind. Each missing free point changes the energy bound by a bounded amount. The total change is linear
in $m$, and its slope $c=c_2+\sigma$ enters $\rho$.

Two terms from Section~\ref{sec:real} are added. The factor $\prod_l(1+\zeta s_l)^m$ of Lemma~\ref{lem:binomvand}
adds the field $\alpha g(s)$, with $g(s)=\log(1+\frac34s)$, to the column field. Since $g>0$ on $(0,1]$, the
value with $\alpha g$ is at most the value with $\gamma g$ for every $\gamma\ge\alpha$. The factor
$e^{N\hat\delta}$ of the weights adds $\hat\delta(\mathsf t(x))/58$ to the field of every prescribed point $x$,
where $\mathsf t(x)=2x/(1+x^2)$. Indeed $n^2\cdot\frac1n\sum_j\hat\delta(\mathsf t(y_j))/58=N\sum_j\hat\delta(\mathsf t(y_j))$.

Every step of OpenAI's argument remains valid uniformly in the number of free points. The regularization of
the logarithmic interactions, the diagonal corrections and the completion of squares only involve sums over
the points present. For the prescribed points of the first kind the relevant field is
\[
 F_1(x)=-\gamma_d\log(1+x^2)-2\mathsf T(u,x)+\tfrac1{58}\hat\delta(\mathsf t(x)),\qquad x\in E_W .
\]
For those of the second kind it is
\[
 F_2(y)=-\gamma_d\log(1+y^2)+\frac{n_0-2g_0-1}{n_0}\log y-u_1y+\tfrac1{58}\hat\delta(\mathsf t(y))
\]
for $y\in E_W$ with $0<y<1$, where $u_1$ is the first coefficient of $u$. The coefficient $(n_0-2g_0-1)/n_0$ is a lower bound, valid for all
$N\ge1$, of $(n-1-2g)/n$. The exact coefficient of $-\log(1+x^2)$ in both fields is $(n-g-1)/n=\gamma_d-1/n$. Since
$0\le\log(1+x^2)\le\log2$ for $|x|\le1$, the exact fields exceed $F_1$ and $F_2$ by at most $\log2/n$, which tends to zero. By Lemma~\ref{lem:binomvand} the generalized Vandermonde determinant contributes
the factor $\prod_{i<l}|s_i-s_l|$ and the column field $\alpha g$ to the energy, and the factors $\mathcal B_\zeta$
are assigned to the columns. The energy bound then reduces to a quadratic inequality in $\alpha$ and $\alpha_2$.
The maximum over $\alpha_2$ has a closed form.

\begin{lemma}[energy inequality with prescribed points]\label{lem:energy-ineq}
Let $\ell=\log2$ and let $\mathcal N,\mathcal F,\mathcal R,\mathcal F_2$ be real numbers. Put
\begin{multline*}
 V(\alpha)=\bigl(-(a_0-\tfrac12)+(a_1-1)\alpha-\alpha^2\bigr)\ell+\mathcal N+(1-\alpha)\mathcal F+\alpha\mathcal R\\
 +\frac{(\mathcal F_2+\ell-\mathcal R)^2}{4\ell}+U_0-\tfrac12\ell-c_2\alpha .
\end{multline*}
Then for all real $\alpha$ and all $\alpha_2\ge0$
\begin{multline*}
 \bigl(-(a_0-\tfrac12)+(a_1-1)\alpha-\alpha^2+\alpha_2-\alpha_2^2\bigr)\ell+\mathcal N+(1-\alpha)\mathcal F\\
 +(\alpha-\alpha_2)\mathcal R+\alpha_2\mathcal F_2\le-U_0+\tfrac12\ell+c_2\alpha+V(\alpha) .
\end{multline*}
\end{lemma}

\begin{proof}
The difference of the two sides is $\alpha_2(\mathcal F_2+\ell-\mathcal R)-\ell\alpha_2^2$, and
$tb-\ell t^2\le b^2/(4\ell)$ for all real $t$ and $b$.
\end{proof}

So the energy bound holds at the slope $c=c_2+\sigma$ on every range of $\alpha$ on which $V(\alpha)\le\sigma\alpha$.
The lemma is applied with $\mathcal N$ a bound for $2\|u\|^2+\frac12\|v\|^2$, $\mathcal F$ a bound for $M_X+M_Y$, with the
column field $\gamma g$ included in $Y$, $\mathcal R$ a bound for the supremum of $F_1$ and $\mathcal F_2$ a bound for the
supremum of $F_2$. Keeping $\mathcal R$ and $\mathcal F_2$ apart is much better than one cap for both, because the moved
trial sequences of \S\ref{ssec:trials} lower the first field and raise the second. The analogue for the constrained case has no prescribed
points of the second kind. Its inequality is
\[
 \bigl(-(a_0-\tfrac12)+(a_1-1)\alpha-\alpha^2\bigr)\ell+n_1+(1-\alpha)\mathcal F_1+\alpha\mathcal M_1
 \le-U_0+\tfrac12\ell+c\alpha,
\]
with $n_1$ a bound for the norms of the trial sequences, $\mathcal F_1=M_{X_1}+M_Y$, and $\mathcal M_1$ a bound for
$\lambda_1\bigl(\frac{2g_0}{n_0}-\frac{2x^2}{1+x^2}\bigr)+\hat\delta(\mathsf t(x))/58$ on $E_W$. Since
$1-\frac{2x^2}{1+x^2}=\sqrt{1-\mathsf t(x)^2}$, this bound is a maximum of affine functions of $\lambda_1$ over the
segments of the chord table, with $\hat\delta$ replaced by its larger end value. So it is convex in $\lambda_1$.

\emph{The squares that remain.} Section~\ref{sec:signed} uses what the completion of squares leaves, so we record it for the case with doubled
interaction. Let $\mu_F$, $\mu_M$ and $\mu_S$ be the normalized counting measures of the free points, the prescribed points of the first kind and the
prescribed points of the second kind. Their total masses are $1-\alpha$, $\alpha-\alpha_2$ and $\alpha_2$. Put $W=\mu_F+\mu_M$,
$A_k(\mu)=\int T_k\,d\mu$ and $\langle a,b\rangle=\sum_ka_kb_k/k$, so that $\|a\|^2=\langle a,a\rangle$. Three products carry the interactions
in the variable $x$. The Vandermonde product of the free points comes from the Cauchy determinant. The row determinant gives the Vandermonde product of the
free points and the prescribed points of the first kind, and the Vandermonde product of the points $y$ of the second kind. These last points
enter the row determinant through $1/y$, so it also gives the cross factors $|w-1/y|=(1-wy)/y$, for $w$ a point of $W$.
Since $-\log|x-x'|=\log2+2\sum_kT_k(x)T_k(x')/k$ and $\log(1-wy)\le-wy$, the part of $E$ quadratic in the Chebyshev moments is, apart from the
diagonal terms and the multiples of $\log2$, at most
\[
 -\|A(\mu_F)\|^2-\|A(W)\|^2-\|A(\mu_S)\|^2-A_1(W)A_1(\mu_S).
\]
For any measure $\mu$ we have $-\|A(\mu)\|^2=\|u\|^2-2\langle u,A(\mu)\rangle-\|A(\mu)-u\|^2$. Using this for $\mu_F$ and $W$, and writing
$A_1(W)=u_1+(A_1(W)-u_1)$, we find that the display equals
\[
 2\|u\|^2-2\langle u,A(2\mu_F+\mu_M)\rangle-u_1A_1(\mu_S)-\mathcal Q,
\]
where
\[
 \mathcal Q=\|A(\mu_F)-u\|^2+\|A(W)-u\|^2+\|A(\mu_S)\|^2+(A_1(W)-u_1)A_1(\mu_S).
\]
The first three terms are the term $2\|u\|^2$ of $\mathcal N$ and the terms of $X$, $F_1$ and $F_2$ that contain $u$. So the energy bound obtained from
Lemma~\ref{lem:energy-ineq} improves by $\mathcal Q$. To bound $\mathcal Q$ from below, let $e_1$ be the first
unit vector and put $a=A(\mu_F)-u$ and $b=A(W)-u+\frac12A_1(\mu_S)e_1$. Since $\|e_1\|=1$,
\[
 \|b\|^2=\|A(W)-u\|^2+(A_1(W)-u_1)A_1(\mu_S)+\tfrac14A_1(\mu_S)^2 .
\]
Also $\frac14A_1(\mu_S)^2\le\|A(\mu_S)\|^2$, so $\mathcal Q\ge\|a\|^2+\|b\|^2\ge2\|\frac12(a+b)\|^2$. The average of $a$ and $b$ is
$A(\mu_F+\frac12\mu_M)-u+\frac14A_1(\mu_S)e_1$. This gives
\begin{equation}\label{eq:squares}
 \mathcal Q\ge\Delta,\qquad\Delta=2\sum_k\frac1k\bigl(A_k(\mu_F+\tfrac12\mu_M)+\tfrac14A_1(\mu_S)\delta_{k1}-u_k\bigr)^2 .
\end{equation}
The weight $\frac12$ of $\mu_M$ records that these points occur in one of the two Vandermonde products of the $x$ variables, and the weight
$\frac14$ is half of the $\frac12$ in $b$. The diagonal terms are treated by OpenAI's regularization, which replaces the moments by damped moments
with an error that tends to zero, as in the proof of Lemma~\ref{lem:freetangent}. We suppress this.

\subsection{The trial sequences}\label{ssec:trials}
The trial sequences have the form
\begin{equation}\label{eq:trial}
 u_k=c_k+\operatorname{Re}\sum_jr_jz_j^k\qquad(k\ge1;\ c_k=0\text{ for }k>\Kcoef),
\end{equation}
with $\Kcoef$ Chebyshev coefficients and finitely many geometric terms. The geometric terms have a simple
meaning. Since $\sum_{k\ge1}z^kT_k(x)/k=-\frac12\log(1-2xz+z^2)$, a term $(z,r)$ contributes
$-\frac12\operatorname{Re}\,r\log(1-2xz+z^2)$ to the potential $\mathsf T(u,x)$. If
$x_0+i\epsilon_0=\frac12(z+z^{-1})$, then $1-2xz+z^2=2z(x_0+i\epsilon_0-x)$. So the term is, up to a constant,
the logarithmic potential of a point charge of complex strength $r$ at $x_0+i\epsilon_0$ and its conjugate. The
extremal measure has a gap in its support around $x=0$. Its potential has a logarithmic singularity at $x=0$
and square-root behaviour at the edges of the gap. A Chebyshev expansion converges slowly near such points,
while a charge at a small distance $\epsilon_0$ reproduces the local behaviour with one term. This is also why
the quadrature nodes of Section~\ref{sec:real} avoid a neighbourhood of $x=0$.

\emph{Moving trial sequences.} The best trial sequence depends on $\alpha$. We use families
\[
 u(\tau)=u+\tau\,\mathrm du,\qquad v(\tau)=v+\tau\,\mathrm dv\qquad(0\le\tau\le1),
\]
with fixed finitely supported increments and with $\tau=\tau(\alpha)$ an explicit piecewise linear function. The fields
$X$, $Y$, $F_1$ and $F_2$ are affine in $\tau$. So the supremum of a field at $\tau$ is at most the chord between the
suprema at $\tau=0$ and $\tau=1$, or, for the kept base fields, the supremum at $\tau=0$ plus $\tau$ times the supremum of the
increment. The norms are quadratic in $\tau$. All bounds then become polynomials in $\alpha$ on each piece of
$\tau(\alpha)$. Appendix~\ref{app:field} gives the charges, the increments and the certified bounds.

\subsection{The two cases at the improved rate}\label{ssec:cases}
Theorem~\ref{thm:real} asks for the rate $U_0=U+\Gamma_0/n_0^2$ instead of $U$, at the slope $c$. In both cases this
reduces to finitely many polynomial inequalities in $\alpha$, each of the form $V(\alpha)\le\sigma\alpha$ on a closed
interval with polynomial caps. Each is proved by a cover of the interval by subintervals on which the leading
term dominates, in exact rational arithmetic. The bound at an interior $\alpha$ of a cell uses the caps at the right
end point of the cell, since $g>0$. A finite cover by cells and the uniform continuity of the caps then give the
bound for all $\alpha$, with an error that tends to zero.

In the constrained case we use the increment with $24$ Chebyshev coefficients and $\tau=\min(1,\frac{50}{11}\alpha)$. The
inequality holds on $[0,1]$, which is covered by the $6$ pieces with end points $0,\frac{11}{50},\frac3{10},\frac25,\frac35,\frac45,1$.
In the case with doubled interaction we use the increment with $40$ Chebyshev coefficients and $\tau=1$. The inequality
of Lemma~\ref{lem:energy-ineq} holds on $[\alpha_1,1]$, covered by the $6$ pieces with end points
$\alpha_1,\frac15,\frac3{10},\frac25,\frac35,\frac45,1$. For $\alpha<\alpha_1$ it fails, because the extra term
$U_0-U$ is absorbed by the term $c\alpha$ only once $\alpha$ is not small. So a new argument is needed only in the case
with doubled interaction with fewer than $\alpha_1n$ prescribed columns. Section~\ref{sec:signed} supplies it, with the
same slope $c$. Precisely, it proves the following.

\begin{proposition}\label{prop:signed-goal}
In the case with doubled interaction, for every configuration with $m\le\alpha_1n$ prescribed columns at the nodes
$t_{a_1},\ldots,t_{a_m}$, the integrand of Lemma~\ref{lem:andreief}, with the generalized Vandermonde determinant replaced by
the bound of Lemma~\ref{lem:binomvand}, is at most
\[
 \exp\Bigl(n^2\bigl(-U_0+\tfrac12\log2+c\alpha+o(1)\bigr)-N\sum_j\hat\delta(t_{a_j})\Bigr)\,C_*^m\prod_j\mathcal B_\zeta(k_j),\qquad \alpha=m/n,
\]
with an absolute constant $C_*$.
\end{proposition}

\begin{proof}[Proof of Theorem~\ref{thm:real}, given Proposition~\ref{prop:signed-goal}]
By Proposition~\ref{prop:signed-goal} and the discussion above, in both cases a determinant $M_f$ with $m$
prescribed columns $k_1,\ldots,k_m$ at the nodes $t_{a_j}$ is at most
\[
 C_*^m\prod_j\mathcal B_\zeta(k_j)\,e^{n_0cNm-N\sum_j\hat\delta(t_{a_j})}\exp\bigl(n_0^2N^2(-U_0+\tfrac12\log2+o(1))\bigr).
\]
The factor $C_*^m$ comes from the expansion of the row functions at the prescribed points. By \eqref{eq:wk} and
\S\ref{ssec:weights}, the weight of the column $k_j$ at the node $t_{a_j}$, times $\mathcal B_\zeta(k_j)$, is at most
$e^{(\psi+\hat\delta(t_{a_j})+o(1))N}$. The terms $\hat\delta$ cancel. So for every column $k$ and every node $a$
the weight and the factor of $M_f$ that belong to the pair $(k,a)$ have the product
\[
 \log(B_kC_k)\le\bigl(n_0c+\psi+o(1)\bigr)N=(\rho+o(1))N ,
\]
and Lemma~\ref{lem:expansion} applies with this bound for each pair.
For $|\varepsilon|\le e^{-\rho N}$ Lemma~\ref{lem:expansion} gives a factor $(1+O(N^5))^n=e^{o(N^2)}$. With
$D=n_0^2(U_0-\frac12\log2)$ this is the claim.
\end{proof}
